\documentclass[11pt]{article}
\usepackage{amsmath,amsthm,amssymb}
\usepackage[margin=1in]{geometry}
\usepackage{booktabs}
\usepackage{array}
\usepackage[colorlinks=true,linkcolor=blue,citecolor=blue]{hyperref}

\newtheorem{theorem}{Theorem}
\newtheorem{proposition}[theorem]{Proposition}
\newtheorem{corollary}[theorem]{Corollary}
\newtheorem{lemma}[theorem]{Lemma}
\theoremstyle{definition}
\newtheorem{definition}[theorem]{Definition}
\newtheorem{remark}[theorem]{Remark}
\newtheorem{problem}[theorem]{Problem}

\newcommand{\Sb}{\mathrm{Sb}}
\newcommand{\SKB}{\mathsf{SKB}}
\newcommand{\Grp}{\mathsf{Grp}}
\newcommand{\Ab}{\mathsf{Ab}}
\newcommand{\RadRng}{\mathsf{RadRng}}
\newcommand{\BR}{\mathsf{BR}}
\newcommand{\Z}{\mathbb{Z}}
\newcommand{\id}{\mathrm{id}}
\newcommand{\im}{\operatorname{im}}
\newcommand{\Aut}{\operatorname{Aut}}
\newcommand{\Soc}{\operatorname{Soc}}
\newcommand{\shat}{\widehat{\sigma}}
\newcommand{\Hgp}{H^2_{\Grp}}

\title{Two legs of degree two:\\
\large the extension-classifying cocycle $H^2$ and the semi-abelian homology $H_2$
of skew braces, and the open bridge between them}
\author{MacBeth\\ \small Kodamai}
\date{30 September 2026}

\begin{document}
\maketitle

\begin{abstract}
Two ``second (co)homologies'' now circulate in skew-brace theory, and it is tempting to
read one as a duplicate --- or a scoop --- of the other. They are not the same invariant
seen twice. The first is a \emph{cohomology} $H^2$ that classifies extensions of a skew
brace by an ideal: its coefficient is a \emph{module}, and its data include a
\emph{residue-dependent} twist coming from the multiplicative action. The second is a
\emph{homology} $H_2$ of Barr--Beck/Hopf-formula type in the semi-abelian category of skew
braces: its coefficient is a choice of \emph{reflector}, its action is \emph{fixed}, and it
is a Baer invariant of a single object rather than a classifier of extensions. We place the
two side by side, catalogue the structural differences, and locate precisely the one place
they should meet: an open universal-coefficients bridge --- ``one and a half hops'' --- from
the homology leg to the cohomology leg. We state that bridge as a problem, not a theorem,
and identify the radical-ring reflector as the natural candidate coefficient for a homology
home of the extension-classifying obstruction.
\end{abstract}

\section{Introduction}\label{sec:intro}

\paragraph{The problem: ``degree two'' is ambiguous.}
Skew braces are the algebraic backbone of set-theoretic solutions to the Yang--Baxter
equation, and their extension theory has become active on two fronts at once. Both fronts
produce something a reader will call a ``second (co)homology of a skew brace,'' and the two
objects are easy to conflate --- they carry the same numeral, they live over the same class
of algebras, and they were both worked out in the same eighteen months.

On the first front, Rathee and Yadav~\cite{RY} classify skew-brace extensions
$0\to D\to G\to H\to 0$ by a second cohomology $H^2_{\Sb}(H,D)$ whose cocycles are
\emph{pairs} $(\beta,\tau)$ --- an additive $2$-cocycle and a multiplicative $2$-cocycle
coupled by a compatibility identity. On the second front, Gran, Letourmy and
Vendramin~\cite{GLV} develop a \emph{homology} of skew braces of Barr--Beck type, presented
by relative Hopf formulae in the semi-abelian category $\SKB$ and yielding a
Stallings--Stammbach five-term exact sequence.

\paragraph{The gap.}
No account has yet said plainly how these two degree-two invariants relate. The silence
invites two opposite errors. One is to assume they coincide --- that the homology $H_2$
\emph{is} the cohomology $H^2$ dualised, so that either program subsumes the other. The
other is to assume they are simply unrelated. Both errors matter here concretely: the
present author's obstruction-theoretic results on skew-brace decomposition --- the bilinear
class $[\Omega]$ and the residue-dependent solvability obstruction $o_\nu$ --- live squarely
on the cohomology leg, and it is a genuine question whether the homology leg scoops them,
duplicates them, or neighbours them.

\paragraph{The result of this note.}
It neighbours them. We give a side-by-side dictionary (Table~\ref{tab:legs}) of the two legs
and argue that they differ in three structural invariants that no change of bookkeeping can
reconcile:
\begin{itemize}
\item[(i)] \textbf{Coefficients.} The cohomology leg takes a \emph{module} $D$; the homology
leg takes a \emph{reflector} $X\in\{\Grp,\RadRng,\BR,\Ab\}$ into a subvariety.
\item[(ii)] \textbf{Residue.} The cohomology leg is \emph{residue-dependent}: its data
include the multiplicative action $\nu$ and the $\circ$-conjugation twist. The homology leg
is \emph{fixed-action}: the actions factor through the abelianisation and there is no residue
datum at all.
\item[(iii)] \textbf{What is classified.} The cohomology leg \emph{classifies extensions};
the homology leg is a \emph{Baer invariant of one object}, and it treats non-abelian kernels
only for \emph{central} extensions.
\end{itemize}
Having separated them, we then locate the single place they \emph{should} meet: a universal-coefficients
passage from $H_2$ to $H^2$, scaffolded by the five-term sequence of~\cite{GLV} but left open
there. We state it as Problem~\ref{prob:bridge} and no more --- this is a positioning note,
not a claim.

\paragraph{Why it matters.}
For the compositional-correctness programme the payoff is a clean map of the territory: the
$[\Omega]$/$o_\nu$ package is not scooped by the semi-abelian homology, because that homology
computes a different invariant of a different flavour; and the two nonetheless share a precise,
buildable target. Naming that target is the point.

\paragraph{Outline.}
Section~\ref{sec:prelim} fixes the minimal vocabulary. Section~\ref{sec:cohom} is the
cohomology leg (Rathee--Yadav, and the author's $[\Omega]$ and $o_\nu$).
Section~\ref{sec:hom} is the homology leg (Gran--Letourmy--Vendramin).
Section~\ref{sec:dictionary} is the dictionary. Section~\ref{sec:bridge} states the open
bridge. Section~\ref{sec:conclusion} draws the consequence for the programme.

\medskip
\noindent\emph{A typographic warning.} Throughout, $H^2$ (superscript) is always the
\emph{cohomology} of Section~\ref{sec:cohom} and $H_2$ (subscript) always the \emph{homology}
of Section~\ref{sec:hom}. The distinction between the two legs is exactly the distinction
between the two positions of the ``$2$''.

\section{Preliminaries}\label{sec:prelim}

A \emph{skew brace} is a set $A$ with two group structures $(A,+)$ and $(A,\circ)$, sharing
a unit $0$, related by the compatibility
\[
a\circ(b+c) = (a\circ b) - a + (a\circ c)\qquad(a,b,c\in A).
\]
Writing $\lambda_a(b) := -a + (a\circ b)$ gives the \emph{$\lambda$-action}
$\lambda\colon (A,\circ)\to\Aut(A,+)$. An \emph{ideal} $D\trianglelefteq A$ is a normal
subgroup of both $(A,+)$ and $(A,\circ)$ that is stable under $\lambda$; it is exactly the
kernel of a skew-brace morphism, so ideals index quotients. A skew brace is a \emph{trivial
brace} when $\circ=+$ (equivalently $\lambda\equiv\id$). The \emph{socle} is
$\Soc(A) = \ker\lambda\cap Z(A,+)$. Skew braces and their morphisms form a variety $\SKB$,
which is \emph{semi-abelian}~\cite{GLV}.

Fix a skew-brace \emph{extension} of a base $H$ by an ideal $D$,
\begin{equation}\label{eq:ext}
0 \longrightarrow D \longrightarrow G \longrightarrow H \longrightarrow 0 ,
\end{equation}
and a normalised set-theoretic section $s\colon H\to G$ with $s(0)=0$. Following~\cite{RY}
one reads off the \emph{good triplet}
\[
\nu\colon (H,\circ)\to\Aut(D,+),\qquad
\mu\colon (H,+)\to\Aut(D,+),\qquad
\sigma\colon (H,\circ)\to\Aut(D,+),
\]
by $\nu_h=\lambda_{s(h)}|_D$, $\mu_h(x)=-s(h)+x+s(h)$, and $\sigma_h(x)=s(h)^{-1}\circ x\circ s(h)$
(the last is the $\circ$-conjugation, which we also write $\shat_h$). The triplet is
independent of the choice of section. We call $\nu$ the \emph{residue} of the
extension; it is the datum by which the multiplicative structure of $G$ acts on the additive
structure of $D$, and it is precisely what the homology leg will not see.

\section{The cohomology leg: extensions and their obstructions}\label{sec:cohom}

\subsection{The Rathee--Yadav classification}
The extensions~\eqref{eq:ext} realising a fixed good triplet are classified by a second
cohomology $H^2_{\Sb}(H,D)$~\cite{RY}. A cocycle is a \emph{pair}
\[
(\beta,\tau),\qquad
\beta\colon H\times H\to D \text{ (additive $2$-cocycle)},\quad
\tau\colon H\times H\to D \text{ (multiplicative $2$-cocycle)},
\]
subject to a compatibility identity coupling $\beta$ and $\tau$ through the actions
$\lambda,\nu,\mu$ (Rathee--Yadav's equation~(3.10)). The class $[(\beta,\tau)]$ vanishes iff
the extension splits, i.e.\ iff $G$ decomposes as an internal Zappa--Sz\'ep-type product of
$H$ and $D$. The coefficient is the \emph{module} $D$ together with its actions; the theory
in~\cite{RY} is developed for $D$ a \emph{trivial brace} kernel.

The point to hold onto is that a single degree-two class already \emph{carries two group
Schreier classes inside it}: forgetting $\tau$ leaves an additive class, forgetting $\beta$
leaves a multiplicative one, and the compatibility identity binds them.

\subsection{The bilinear obstruction $[\Omega]$}
For the decomposition problem ``does $G$ split as $(G/D)\bowtie D$ along an ideal $D$?'' the
obstruction is a class $[\Omega]$, and it \emph{is} the Rathee--Yadav class $[(\beta,\tau)]$
in the trivial-kernel case~\cite{RY}; this identification has been peer-reviewed within the
group. Its structure is genuinely bilinear, and the following facts pin down how:
\begin{itemize}
\item The \emph{additive forgetful map} $\varphi_+\colon [(\beta,\tau)]\mapsto[\beta]$ into
the group cohomology $\Hgp((H,+),D;\mu)$ is well defined, with no extra hypothesis: a change
of section shifts $\beta$ by the standard additive $2$-coboundary, independently of $\tau$.
(The \emph{multiplicative} forgetful map $\varphi_\circ\colon[(\beta,\tau)]\mapsto[\tau]$ is
Rathee--Yadav's own.)
\item The \emph{product map} $\Phi=(\varphi_\circ,\varphi_+)$ into
$H^2_{\Grp}((H,\circ),D)\times \Hgp((H,+),D;\mu)$ \emph{fails to be injective}: on the
$4\times4$ witness $H=\Z/2\times\Z/4$ there is a class with $[\Omega]\neq0$ but
$\Phi([\Omega])=(0,0)$. The obstruction is therefore irreducibly bilinear --- it lives in the
coupling and is strictly stronger than either group Schreier class.
\end{itemize}
Both the descent of $\varphi_+$ and the $\Phi$-non-injectivity are proved, and the crux lemma
(the section-change computation) is machine-checked in Lean~4.

\subsection{The residue-dependent solvability obstruction $o_\nu$}
The forgetful maps raise the reverse question: which additive classes
$[\beta]\in \Hgp((H,+),D;\mu)$ actually \emph{lift} to a full skew-brace cocycle inducing the
fixed triplet --- that is, what is $\im\varphi_+$? The answer is an obstruction $o_\nu$, and
its defining feature is that it is \emph{residue-dependent}: the additive class $[\beta]$, the
residue $\nu$, and the multiplicative twist are \emph{coupled}, not independent.

\begin{theorem}[trivial-brace-kernel regime; proved]\label{thm:onu}
In the Rathee--Yadav (trivial-brace-kernel) setting, $o_\nu$ is a well-defined group
homomorphism
\[
o_\nu\colon \Hgp((H,+),D;\mu)\longrightarrow C^3\big/A(Z^2_\circ),
\qquad [\beta]\mapsto[Q\beta],
\]
killing the additive coboundaries, and
\[
\im\varphi_+ \;=\; \ker o_\nu .
\]
Equivalently, an additive class lifts to a skew-brace cocycle with the given triplet iff
$o_\nu([\beta])=0$.
\end{theorem}

The proof is structural: the coupling operators are $\Z$-linear, a ``linchpin'' coboundary
identity shows $o_\nu$ kills $B^2_+$, and Rathee--Yadav's cocycle characterisation supplies
the surjection onto genuine cocycles as a cited black box. The $|G|=8$ collapsed slice
($H=\Z/2$, $D=\Z/4$) is formalised sorry-free in Lean~4. Two honesty flags travel with the
theorem: (a) it is proved \emph{in the trivial-brace-kernel regime}, resting on~\cite{RY}'s
characterisation as a premise; and (b) beyond that regime the picture is sharper but not yet
fully proved, as follows.

\begin{remark}[non-trivial kernel; proved at $H=\Z/2$, computed in general]\label{rmk:nontriv}
When $D$ is a \emph{non-trivial} brace, the correct home of the multiplicative datum is the
group $(D,\circ)$ twisted by the $\circ$-conjugation $\shat$ --- not the additive $(D,+)$ ---
and there is an intrinsic \emph{defect} $\delta_h=\nu_h\shat_h-\id$ measuring the failure of
the trivial-kernel collapse. For $H=\Z/2$ this yields a proved ``three-way lock''
$(\lambda^D_x-1)\,b = 2(\nu_1\shat_1(x)-x)$ (verified on all $112$ genuine kernels), and
$\im\varphi_+=o_\nu^{-1}(0)$ becomes an \emph{affine} coset: it is a torsor with no basepoint
exactly when $L:=\im\lambda^D\subsetneq\Aut(D)$ and $\nu\shat$ is non-socle. The general-$H$
form of the obstruction is validated at operator level ($16/16$ triplets for $H=\Z/2$, exact
assembly $33/33$) but its closed-form offset for general $H$ is \emph{computed}, not yet
proved --- one symbolic step remains open.
\end{remark}

\subsection{The character of the leg}
Three features recur. The coefficient is a \emph{module} $D$. The theory is
\emph{residue-dependent}: $\nu$ and the $\circ$-twist are part of the data, and $o_\nu$ exists
precisely to measure their coupling to $[\beta]$. And the invariant \emph{classifies
extensions}: a class is an extension up to equivalence, and it vanishes iff the extension
splits. Every one of these three is reversed on the other leg.

\section{The homology leg: Baer invariants of one object}\label{sec:hom}

\subsection{Barr--Beck homology via Hopf formulae}
Gran, Letourmy and Vendramin~\cite{GLV} compute a \emph{homology} of skew braces in the
semi-abelian category $\SKB$. It is comonadic (Barr--Beck) homology, made explicit by a
\emph{relative Hopf formula}: for a free presentation $0\to K\to F\to B\to0$ and a reflector
$F_X\colon\SKB\to X$ onto a subvariety $X$,
\begin{equation}\label{eq:hopf}
H_2(B,\,F_X) \;\cong\; \frac{K\cap[F,F]_X}{[K,F]_X},
\qquad
H_1(B,\,F_X) \;\cong\; \frac{B}{[B,B]_X},
\end{equation}
where $[\,\cdot\,,\,\cdot\,]_X$ is the $X$-relative commutator. The genuinely new algebraic
object is the \emph{radicalator} $[I,A]_{\RadRng}$, the relative commutator for the reflector
onto radical rings; the Huq and Smith commutators coincide with the abelian commutator
$[\,\cdot\,,\,\cdot\,]_{\Ab}$. The main extension-level output is a Stallings--Stammbach
\emph{five-term exact sequence} for a short exact sequence of skew braces.

\subsection{The character of the leg}
The coefficient here is not a module: it is a \emph{choice of reflector}
$X\in\{\Grp,\RadRng,\BR,\Ab\}$, i.e.\ a choice of ``abelianisation direction'' into a
subvariety. The theory is \emph{fixed-action}: the $\lambda$- and $\rho$-actions factor
through the additive abelianisation $A_+/J_+$, so there is \emph{no residue datum} $\nu$ ---
the very thing the cohomology leg is built around is quotiented away here by construction.
And the invariant is a \emph{Baer invariant of one object} $B$: $H_2(B,F_X)$ is an invariant
of $B$ (and $X$), computed from any free presentation, not a set of extensions of $B$ by a
prescribed module. Non-abelian kernels are treated only for \emph{central} extensions, and
the paper carries no numerical examples --- a size-$24$ database counterexample is invoked
only to \emph{separate} commutators, not to compute a class.

\section{The dictionary}\label{sec:dictionary}

Table~\ref{tab:legs} is the centre of this note. Read down each column and the two legs are
seen to disagree on every structural axis; read across each row and one sees why neither can
be a relabelling of the other.

\begin{table}[h]
\centering
\renewcommand{\arraystretch}{1.35}
\begin{tabular}{@{}p{0.26\textwidth} p{0.34\textwidth} p{0.34\textwidth}@{}}
\toprule
& \textbf{Cohomology leg} (\S\ref{sec:cohom}) & \textbf{Homology leg} (\S\ref{sec:hom})\\
\midrule
Invariant & $H^2_{\Sb}(H,D)$ (superscript) & $H_2(B,F_X)$ (subscript)\\
Type & extension-\emph{classifying} & \emph{Baer invariant} of one object\\
Coefficient & a \emph{module} $D$ (with actions) & a \emph{reflector} $X\in\{\Grp,\RadRng,\BR,\Ab\}$\\
Cocycle / formula & pair $(\beta,\tau)$ + compatibility (3.10) & Hopf formula $(K\cap[F,F]_X)/[K,F]_X$\\
Action & \emph{residue-dependent}: $\nu$, $\shat$ are data & \emph{fixed}: actions factor through $A_+/J_+$\\
Residue $\nu$? & yes --- $o_\nu$ measures its coupling & no --- quotiented away\\
Non-abelian kernel & general ideal $D$ & \emph{central} extensions only\\
Vanishing means & extension splits & (leg-specific; not a splitting criterion)\\
Extension output & the classification itself & Stallings--Stammbach $5$-term sequence\\
Examples & explicit ($\Z/2\times\Z/4$, $\Z/8$, \dots) & none (separating counterexample only)\\
Author's stake & $[\Omega]$, $\varphi_\pm$, $o_\nu$ live here & (neighbour; no overlap)\\
Sources & Rathee--Yadav~\cite{RY} & Gran--Letourmy--Vendramin~\cite{GLV}\\
\bottomrule
\end{tabular}
\caption{The two legs of degree two. The single class $H^2_{\Sb}$ classifies extensions with
a residue-dependent module coefficient; the homology $H_2$ is a fixed-action Baer invariant
with a reflector coefficient. No column agrees.}
\label{tab:legs}
\end{table}

The decisive rows are \emph{coefficient}, \emph{residue}, and \emph{type}. A reflector is not
a module and a module is not a reflector; a residue-dependent theory cannot be a
fixed-action one; and a classifier of extensions is not a Baer invariant of a single object.
These are not bookkeeping differences that a dualisation would erase.

\section{The open bridge}\label{sec:bridge}

And yet the two legs are not strangers. Degree-two homology and degree-two cohomology are
related, in ordinary algebra, by universal coefficients: $H^2(-;M)$ is built from $H_2(-)$ by
dualising into $M$, modulo an $\mathrm{Ext}$ term. The question is whether the semi-abelian
$H_2$ of~\cite{GLV} feeds the extension-classifying $H^2$ of~\cite{RY} the same way. The
honest description of the distance is \emph{one and a half hops}:

\begin{enumerate}
\item[\textbf{Hop 1.}] From the homology $H_2(B,F_X)$ to a derived functor of abelianisation.
This is exactly what the Hopf formula~\eqref{eq:hopf} provides, and~\cite{GLV}'s five-term
sequence is the scaffold linking $H_2$, $H_1$, and the abelianised kernel across an extension.
\item[\textbf{Half a hop.}] From that derived functor, \emph{dualise / apply universal
coefficients} into a module coefficient $D$ to reach a cocycle $H^2$ with module
coefficients. This is the step that is \emph{not} carried out in~\cite{GLV} --- it is left
open there --- and it is only ``half'' a hop because the target is not the na\"\i ve
$H^2(-;D)$ but the residue-carrying $H^2_{\Sb}(H,D)$, and the residue $\nu$ must be
\emph{reintroduced} on the cohomology side after the homology side has quotiented it away.
\end{enumerate}

\begin{problem}[the bridge]\label{prob:bridge}
Construct a universal-coefficients passage from the semi-abelian homology $H_2(B,F_X)$
of~\cite{GLV} to the extension-classifying cohomology $H^2_{\Sb}(H,D)$ of~\cite{RY}, with the
five-term sequence of~\cite{GLV} as scaffold. In particular, determine which reflector $X$ and
which coefficient recover the residue-dependent data $(\nu,\shat)$ of the cohomology leg, so
that the obstructions $[\Omega]$ and $o_\nu$ acquire a homological description.
\end{problem}

We do not claim this bridge; we state it. Two remarks constrain the search. First, because the
homology leg is fixed-action, any bridge must \emph{reintroduce} the residue --- so a plain
dualisation cannot suffice, and the interesting content is exactly the recovery of $\nu$.
Second, the radical-ring reflector is the natural candidate coefficient: the radicalator
$[I,A]_{\RadRng}$ and $\RadRng$-centrality are the part of~\cite{GLV} that sees multiplicative
(as opposed to purely additive) structure, which is where a homological home for the
bilinear, $\circ$-twisted obstruction $[\Omega]$ would have to live. We flag
$X=\RadRng$ as the candidate $[\Omega]$-analogue home and leave the construction open.

\section{Consequence for the programme}\label{sec:conclusion}

The extension-classifying obstructions $[\Omega]$ and $o_\nu$ have a published homology
\emph{neighbour} --- the semi-abelian $H_2$ of Gran--Letourmy--Vendramin --- which they
neither scoop nor duplicate. The two compute different invariants of different type over the
same algebras: a residue-dependent, module-coefficient classifier of extensions on one side,
a fixed-action, reflector-coefficient Baer invariant on the other. For the
compositional-correctness programme this is the clean statement we wanted: the cohomological
work is not pre-empted by the homological work, and the honest joint target is the bridge of
Problem~\ref{prob:bridge}, whose natural coefficient we conjecture to be the radical-ring
reflector. Building that bridge --- reintroducing the residue on the homology side --- is the
next question, and it is a question about both legs at once.

\medskip
\noindent\emph{Provenance and scope.} The Rathee--Yadav classification and their
multiplicative forgetful map are cited, not reproved. The identification $[\Omega]=[(\beta,\tau)]$
is peer-reviewed within the group; the descent of $\varphi_+$ and the non-injectivity of
$\Phi$ are proved and Lean-checked. Theorem~\ref{thm:onu} is proved in the trivial-brace-kernel
regime, resting on~\cite{RY}'s cocycle characterisation as a premise; the non-trivial-kernel
sharpening of Remark~\ref{rmk:nontriv} is proved at $H=\Z/2$ and computed in general, with one
symbolic step open. The facts attributed to~\cite{GLV} are from a full-text reading of that
paper. The bridge of Section~\ref{sec:bridge} is open and is not claimed. Gran--Letourmy--Vendramin
themselves cite Rathee--Yadav's earlier 2024 cohomology paper (a distinct paper from the
extension paper~\cite{RY} used here), a modest pre-existing bibliographic contact point between
the semi-abelian and skew-brace traditions --- but not the bridge, which remains to be built.

\begin{thebibliography}{9}
\bibitem{RY}
M.~Rathee and M.~K.~Yadav,
\emph{Skew brace extensions, second cohomology and complements},
arXiv:2601.12371.
\bibitem{GLV}
M.~Gran, T.~Letourmy and L.~Vendramin,
\emph{Hopf formulae for the homology of skew braces},
J. Pure Appl. Algebra \textbf{229} (2025), no.~11, 108085; arXiv:2409.18056.
\end{thebibliography}

\end{document}
